\subsection{Décomposition du travail sur deux nœuds}
Le calcul distribué compare le coût d'un échange par ligne à celui d'une réduction en fin de lot. Dans le premier parcours d’une paire, chaque rang possède une bande verticale. Au début de la ligne \(i\), il connaît le rang et la colonne de la dernière occurrence antérieure de \(A_i\). Le bord droit de la ligne \(i-1\) permet aussi de traiter le cas où cette occurrence se trouve en première colonne d'une bande.
L'algorithme~\ref{alg:mpi} décrit les deux scalaires dont chaque rang a besoin pour poursuivre sa bande locale.
\begin{algorithm}[!htbp]
\caption{Décomposition MPI d'une paire par bandes de colonnes}\label{alg:mpi}
\begin{algorithmic}[1]
\Require Strings \(A,B\), \(r\) MPI ranks and one column interval per rank
\Ensure The terminal score \(d(A,B)\) on every rank
\State Index the last occurrence of each symbol in each column interval
\For{each row \(i=1,\ldots,m\)}
  \State Exchange each previous right edge and the score before the last local match using \texttt{MPI\_Allgather}
  \State Update the local interval using Algorithm~\ref{alg:nvpd} and swap score rows
\EndFor
\State Broadcast the terminal score from its owning rank
\end{algorithmic}
\end{algorithm}
\schema{mpi}{Le parcours collectif échange deux scores par rang à chaque ligne. Une réduction rassemble les scores du lot.}{fig:comm}{src/mpi/mpi.cpp}
Avec le collectif et \(r\) rangs, chacun reçoit \(2r\) entiers par ligne. Dans le lot de la figure~\ref{fig:comm}, le rang \(q\) prend les paires d'indices \(q,q+r,\ldots\). Chaque fil traite une paire en SIMD. Le temps mesuré est le maximum des temps de rang, hors lecture du FASTA.
\chart{distributed/echelle_mpi.png}{À charge fixe, le rapport des temps sur un et deux nœuds distingue le ralentissement des paires isolées du gain obtenu sur les lots.}{fig:mpifort}{bench/mpi/scaling.py}
La figure~\ref{fig:mpifort} donne la moyenne de sept blocs indépendants après deux échauffements. L’ordre des essais est randomisé. Chaque configuration réserve un ou deux nœuds entiers, soit 192 ou 384 cœurs. Sur la paire de 20\,000 lignes, le temps passe de \(91{,}49\pm0{,}78\) à \(997{,}1\pm12{,}5\)\,ms. Le temps du lot synthétique de 2\,016 paires passe de \(768{,}17\pm3{,}62\) à \(528{,}33\pm0{,}93\)\,ms.

\synthese{Sur deux nœuds, l'efficacité forte atteint 72,70\,\% sur le lot synthétique et 71,35\,\% sur le phage. Elle tombe à 4,59\,\% sur la paire carrée et à 13,70\,\% sur la paire large.}

\subsection{Temps passé dans les échanges collectifs}\label{sec:collectives}
L'interface de profilage PMPI, définie par le MPI Forum, permet d'entourer un appel MPI d'un chronomètre. Après deux échauffements, sept blocs comparent le même calcul avec et sans sonde. Le rang critique est celui qui détermine la durée du calcul. Pour chaque essai instrumenté, la part du collectif est le rapport de son temps cumulé dans \texttt{MPI\_Allgather} à sa durée totale. Le temps ainsi attribué au collectif comprend l'attente des autres rangs.

\chart{distributed/collectives_mpi.png}{Temps du collectif sur le rang critique, coût de l'instrumentation et gain du pipeline de voisinage sur les mêmes deux nœuds.}{fig:collectives}{bench/mpi/collectives/probe.c}
Sur la paire carrée de la figure~\ref{fig:collectives}, les 20\,000 collectifs occupent \(70{,}41\pm1{,}35\)\,\% de ce temps sur un nœud et \(97{,}10\pm0{,}10\)\,\% sur deux. Sur la matrice large, les 1\,024 collectifs en occupent respectivement \(26{,}13\pm1{,}11\) et \(87{,}94\pm0{,}16\)\,\%. L'écart-type porte sur les sept rapports. Pour le carré à deux nœuds, activer la sonde change le temps de \(-4{,}95\pm10{,}07\)\,\% par rapport au témoin apparié. La dispersion empêche d'attribuer cette différence à la seule instrumentation.

Le panneau c) de la figure~\ref{fig:collectives} mesure une version qui transmet le bord droit d'une bande au rang voisin par paquets de \(H\) lignes. Dans une bande dont la colonne locale \(x\) correspond à la colonne globale \(s+x\), posons \(P_x=D_{i-1,s+x}\) et \(C_x=D_{i,s+x}\). Le voisin de gauche fournit \(C_0\). Si \(k\) désigne la dernière occurrence locale de \(A_i\) à gauche de \(x\), y compris \(x\), la récurrence devient
\[
 C_x=\min\bigl\{P_x+1,\,P_{x-1}+1,\,C_0+x,\,P_{k-1}+x-k\bigr\}.
\]
Le dernier terme disparaît en l'absence de correspondance locale. Le terme \(C_0+x\) couvre le chemin qui entre dans la ligne par le bord et se poursuit par \(x\) insertions. Avec \(r\) bandes non vides, les rangs échangent \((r-1)\lceil m/H\rceil\) messages de voisinage. La sonde confirme que les mises à jour n'appellent plus \texttt{MPI\_Allgather}.

La comparaison repose sur sept blocs appariés, deux échauffements et le même placement que le témoin. Sur deux nœuds, le carré demande \(996{,}89\pm14{,}98\)\,ms avec le collectif par ligne et \(13{,}074\pm0{,}425\)\,ms avec \(H=32\). Sur la paire large, les temps passent de \(51{,}925\pm0{,}570\) à \(5{,}815\pm0{,}103\)\,ms avec \(H=8\). Ce gain à ressources fixées se distingue du passage d'un nœud à deux. Le rapport de temps du pipeline seul vaut alors 0,871 pour le carré et 1,372 pour la paire large. Sur cette dernière, \(H=32\) ne produit que 32 paquets pour remplir 48 bandes, contre 128 paquets à \(H=8\).

\synthese{À deux nœuds, le gain apparié sur le témoin atteint 76,30 sur le carré (IC à 95\,\% de 74,52 à 78,07) et 8,93 sur la paire large (8,77 à 9,09). La hauteur la plus rapide dépend de la forme du rectangle.}

\subsection{Efficacité à charge par cœur constante}
L'expérience faible compare deux nœuds complets à un nœud complet, avec huit fils par rang. Soit \(W\) le nombre de cases traitées sur le premier nœud. En doublant le travail, l'efficacité faible et le gain de débit valent
\[
 E_{\mathrm{weak}}=\frac{T_{192}(W)}{T_{384}(2W)},\qquad
 \frac{2W/T_{384}(2W)}{W/T_{192}(W)}=2E_{\mathrm{weak}}.
\]
Pour une paire, 1\,024 lignes sont conservées et la largeur vaut 256 ou 4\,096 colonnes par cœur. Pour le lot, les séquences comptent 8\,192 bases. Les 85 puis 120 séquences donnent respectivement \(85\times84/2=3\,570\) et \(120\times119/2=7\,140\) paires. Le travail total double dans les trois cas. Dans le lot, le nombre de cases par cœur est une moyenne sur les travailleurs.
\chart{distributed/mpi_faible.png}{Au doublement du travail et du nombre de cœurs, sept blocs donnent les temps et l’efficacité faible. Les crochets délimitent les intervalles de confiance à 95\,\%.}{fig:mpifaible}{bench/mpi/scaling.py}
Sur la figure~\ref{fig:mpifaible}, le passage à deux nœuds porte la durée du lot de \(691{,}97\pm7{,}28\) à \(890{,}94\pm3{,}47\)\,ms. Sur une paire, l'efficacité passe de 10,07 à 30,57\,\% lorsque le travail local par ligne est multiplié par seize. Le nombre de collectifs reste fixé à 1\,024.

\synthese{Pour le lot indépendant, l'efficacité faible atteint \(77{,}67\pm0{,}76\)\,\%. Sur la paire large, l'IC à 95\,\% s'étend de 19,33 à 41,80\,\%. La dispersion sur un nœud limite la précision du rapport.}

\subsection{Répartition de fenêtres homologues entre nœuds}\label{sec:mpihumain}
Sur le chromosome 21, un locus de GRCh38 est comparé à son homologue dans CHM13. La chaîne de correspondance UCSC de \texttt{hg38} vers \texttt{hs1}~\cite{ucscchain} fournit les deux bornes. Les 41 bases autour de chaque borne sont identiques dans les deux assemblages. Les 768 fenêtres occupent des intervalles disjoints de 16\,384 bases dans GRCh38. Les insertions et les délétions donnent des longueurs de 5\,297 à 19\,840 bases dans CHM13. Les distances vont de 0 à 11\,107 éditions. Edlib et WFA2 ont retrouvé chaque score.

La commande \texttt{pairs} distribue les comparaisons entre loci homologues. Pour le test d'échelle faible, le lot de 384 reprend une fenêtre sur deux de la même grille chromosomique. Si \(m_i,n_i\) sont les longueurs au locus \(i\), le travail effectif du lot \(L\) et l'efficacité en débit par cœur s'écrivent
\[
 W(L)=\sum_{i\in L}m_i n_i,\qquad
 E_{\mathrm{weak}}=\frac{W(L_2)}{2W(L_1)}\frac{T_{192}(L_1)}{T_{384}(L_2)}.
\]
Le rapport \(W(L_2)/W(L_1)\) vaut 1,9982. Pour chaque configuration, chacun des cinq blocs lance ses propres rangs et chronomètre le deuxième appel après un échauffement complet. Le chronomètre commence après la diffusion du FASTA et retient le temps du rang le plus lent.

\chart{distributed/echelle_mpi_humain.png}{Les échanges synchronisent les rangs sur une paire (a). Les lots de fenêtres se partagent en paires indépendantes (b). Les intervalles de confiance portent sur cinq blocs.}{fig:mpihumain}{bench/mpi/human/measure.py}
Dans le panneau a) de la figure~\ref{fig:mpihumain}, les paquets de 128 lignes alignent la paire de 131\,072 bases en \(159{,}52\pm3{,}91\)\,ms sur un nœud puis \(116{,}89\pm1{,}44\)\,ms sur deux. Avec l'échange collectif par ligne, cette dernière durée atteint \(6{,}672\pm0{,}182\)\,s. Dans le panneau b), la distribution du lot complet porte le débit de \(316{,}5\pm9{,}4\) à \(511{,}6\pm29{,}4\)\,GCUPS. Son efficacité forte vaut 80,84\,\% et son efficacité faible corrigée du travail 80,89\,\%. Chaque rang conserve le FASTA complet. Le plus grand pic RSS parmi les rangs passe de \(44{,}96\pm1{,}00\) à \(45{,}12\pm0{,}28\)\,Mio pour les 768 paires.

\synthese{Sur la paire de 131\,072 bases à deux nœuds, les paquets de 128 lignes divisent le temps par 57,09 face à l’échange collectif par ligne. Sur le lot de 768 paires, répartir des alignements entiers donne 80,8\,\% d’efficacité forte.}

